\subsection{Preliminary downstream signal: LLM test generation}
\label{sec:downstream-prelim}

The structural extension reported in Sections~\ref{sec:empirical} is a
prerequisite for downstream impact, not a demonstration of it. As a
preliminary downstream probe we compare two configurations of an
LLM-based test generator: \emph{P3} (the isolated baseline spec is
attached to the class under test as JML) and \emph{P3C} (the
compositionally-refined spec is attached instead). Both feed the
class source plus the spec into \texttt{gemini-2.5-flash} with
temperature $0.3$, $5$ generation runs per class. The 11-class
subset is the same Commons Lang 3 fa\c{c}ade used in the
predecessor LLM study~\cite{edmonds2026inference}.

Table~\ref{tab:downstream} reports the mean number of test methods
the LLM emits per run, for each class. Three classes
(\texttt{NumberUtils}, \texttt{MutableDouble}, \texttt{Range})
produced no extractable code in either phase (the model's response
exceeded our 16{,}384-token output budget before the closing markdown
fence, and the extractor reports zero tests in that case); these
rows are honest zeros, not failures of the comparison.

\begin{table}[ht]
\centering
\small
\caption{Mean number of test methods extracted per generation run on each Commons Lang class. P3: LLM context contains the isolated baseline spec. P3C: LLM context contains the compositionally-refined spec. Eight of eleven classes returned extractable tests in both configurations; six of those eight show P3C $>$ P3. The largest gains are on validation-heavy classes (\texttt{Validate}, \texttt{BooleanUtils}), where the refined spec's branch-conditional clauses provide additional test-shape hints.}
\label{tab:downstream}
\begin{tabular}{lrrr}
\toprule
\textbf{Class} & \textbf{P3} & \textbf{P3C} & \textbf{$\Delta$} \\
\midrule
BooleanUtils  & 40.0 & 51.0 & \textbf{+11.0 (+27.5\%)} \\
CharUtils     & 58.0 & 44.6 & $-13.4$ ($-23.1\%$) \\
MutableInt    & 21.0 & 21.6 & $+0.6$ ($+2.9\%$) \\
MutableBoolean & 43.2 & 44.4 & $+1.2$ ($+2.8\%$) \\
Pair          & 56.0 & 63.4 & $+7.4$ ($+13.2\%$) \\
MutablePair   & 40.6 & 38.6 & $-2.0$ ($-4.9\%$) \\
Validate      & 39.8 & 61.6 & \textbf{+21.8 (+54.8\%)} \\
Mutable       & 24.0 & 26.8 & $+2.8$ ($+11.7\%$) \\
\midrule
\textbf{Mean (8 classes)} & \textbf{40.3} & \textbf{44.0} & \textbf{+3.7 (+9.1\%)} \\
\bottomrule
\end{tabular}
\end{table}

\textbf{Interpretation, with caveats.} The +9.1\% mean test-count
delta is a weak positive signal: the compositionally-refined context
results in the LLM emitting more test methods per run on six of eight
classes, with the largest gains on validation-heavy code where the
branch-conditional clauses give the LLM additional case structure
to test. The single substantial regression is \texttt{CharUtils}
($-23\%$), where the refined spec adds disjunctive obligations from
the class's many overloaded \texttt{is*} methods; the LLM appears
to interpret the disjunction as one combined assertion rather than
several test cases. The aggregate is positive but the per-class
variance is high, and the n=11 sample size precludes a meaningful
significance test.

\textbf{Threats specific to this measurement.} We report extracted
test count, not mutation-killing rate. The generated tests did not
compile cleanly in our environment --- a tooling issue (LLM output
sometimes includes leading markdown fences our extractor did not
strip, and Gemini-2.5 occasionally emits surrogate-pair characters
in \texttt{char} literals) that prevented us from running the
mutation-coverage step before submission. We commit the
generated test files to the replication package so independent
parties can post-process and run PIT~\cite{coles2016pitest} against
them. A full mutation-coverage comparison is the
top-priority follow-up; this subsection establishes only that the
P3 vs P3C delta is non-zero on the test-generation step itself.
